% =========================================================
% Interval Arithmetic
% =========================================================

\section{Interval Arithmetic}

\begin{propositionbox}{Definition (Interval Addition)}
\label{def:interval-add}
For closed intervals $[a,b]$ and $[c,d]$, their \emph{sum} is
\[
  [a,b] + [c,d] \;=\; [a+c,\; b+d].
\]
\end{propositionbox}

\begin{propositionbox}{Definition (Interval Subtraction)}
\label{def:interval-sub}
For closed intervals $[a,b]$ and $[c,d]$, their \emph{difference} is
\[
  [a,b] - [c,d] \;=\; [a-d,\; b-c].
\]
\end{propositionbox}

\begin{propositionbox}{Definition (Interval Multiplication)}
\label{def:interval-mult}
For closed intervals $[a,b]$ and $[c,d]$, their \emph{product} is
\[
  [a,b] \cdot [c,d] \;=\;
  \bigl[\min(ac, ad, bc, bd),\; \max(ac, ad, bc, bd)\bigr].
\]
\end{propositionbox}

\begin{propositionbox}{Definition (Interval Containment)}
\label{def:interval-contain}
$[a,b] \subseteq [c,d]$ if and only if $c \leq a$ and $b \leq d$.
\end{propositionbox}

\begin{remark*}[Motivation]
Interval arithmetic provides a framework for rigorous numerical computation
with guaranteed error bounds. A real number $x \in [a,b]$ means the computed
value may lie anywhere in the interval; operations propagate these bounds.
\end{remark*}

\begin{remark*}[Status]
\textit{This section is a stub. A full treatment -- the sub-distributivity law,
dependency problem, applications to root-finding, and connections to
compactness -- will be developed in a future revision.}
\end{remark*}
